% ==============================================================================
% annexes/tools.tex — UE SIS589
% Annexe A : Arsenal Outillage — Référence Rapide
% ==============================================================================

\chapter{Arsenal Outillage: Référence Rapide}
\label{ann:tools}

\begin{boxinfo}
Cette annexe recense les outils cités dans le cours, classés par domaine.
Pour chaque outil~: nom, usage principal, plateforme, licence et commande
de référence. Toutes les commandes supposent une installation standard~;
se référer à la documentation officielle pour les options avancées.
\end{boxinfo}

% ==============================================================================
\section{Supervision et SIEM}
\label{ann:tools-siem}
% ==============================================================================

\begin{table}[H]
\centering\small
\caption{Outils de supervision et SIEM}
\label{tab:tools-siem}
\rowcolors{2}{TableauPair}{TableauImpair}
\begin{tabular}{p{2.8cm}p{4cm}p{2.5cm}p{2cm}p{2.5cm}}
\toprule
\tableauHeader{Outil} & \tableauHeader{Usage} &
\tableauHeader{Plateforme} & \tableauHeader{Licence} &
\tableauHeader{URL} \\
\midrule
\textbf{Wazuh} & HIDS / SIEM / FIM / Compliance & Linux, Windows & Open Source &
\texttt{wazuh.com} \\
\textbf{Elasticsearch} & Indexation et recherche fulltext & Linux & Apache 2.0 &
\texttt{elastic.co} \\
\textbf{Kibana} & Visualisation et tableaux de bord & Linux & SSPL &
\texttt{elastic.co} \\
\textbf{Logstash} & Collecte et pipeline de logs & Linux & Apache 2.0 &
\texttt{elastic.co} \\
\textbf{Suricata} & NIDS/NIPS réseau & Linux & GPLv2 &
\texttt{suricata.io} \\
\textbf{Snort~3} & NIDS signatures & Linux, Windows & GPLv2 &
\texttt{snort.org} \\
\textbf{Zeek} & Analyse comportementale réseau & Linux & BSD &
\texttt{zeek.org} \\
\textbf{TheHive} & Gestion d'incidents + ticketing & Linux & Apache 2.0 &
\texttt{thehive.io} \\
\textbf{MISP} & Partage de CTI (IoC) & Linux & Apache 2.0 &
\texttt{misp-project.org} \\
\textbf{OpenCTI} & Plateforme CTI graphique & Linux & Apache 2.0 &
\texttt{filigran.io} \\
\textbf{Velociraptor} & EDR / DFIR & Linux, Windows & Apache 2.0 &
\texttt{docs.velociraptor.app} \\
\bottomrule
\end{tabular}
\end{table}

\begin{lstlisting}[style=StyleTerminal,
  caption={Commandes essentielles: Wazuh},
  label={lst:tools-wazuh}]
# État des agents
/var/ossec/bin/agent_control -l

# Redémarrer le manager
systemctl restart wazuh-manager

# Tester une règle (logtest)
/var/ossec/bin/wazuh-logtest -V

# Recompiler les listes CDB
/var/ossec/bin/ossec-makelists

# Voir les alertes en temps réel
tail -f /var/ossec/logs/alerts/alerts.log

# Forcer une analyse FIM immédiate
/var/ossec/bin/agent_control -a -u AGENT_ID
\end{lstlisting}

% ==============================================================================
\section{Collecte de logs}
\label{ann:tools-logs}
% ==============================================================================

\begin{lstlisting}[style=StyleTerminal,
  caption={Commandes essentielles: collecte de logs},
  label={lst:tools-logs}]
# ─── Rsyslog : démarrage et vérification ─────────────────────────────────────
systemctl status rsyslog
rsyslogd -N1            # Tester la configuration
journalctl -u rsyslog -f

# ─── Logstash : démarrage et test pipeline ────────────────────────────────────
systemctl start logstash
/usr/share/logstash/bin/logstash --config.test_and_exit \
  -f /etc/logstash/conf.d/

# ─── Filebeat : état et configuration ────────────────────────────────────────
filebeat test config
filebeat test output
systemctl status filebeat

# ─── NXLog (Windows) ─────────────────────────────────────────────────────────
# Installer depuis https://nxlog.co/
# Fichier config : C:\Program Files\nxlog\conf\nxlog.conf
net start nxlog

# ─── Windows Event Forwarding (WEF) ──────────────────────────────────────────
# Sur le collecteur Windows
wecutil qc       # Configurer le collecteur
wecutil cs subscription.xml  # Créer un abonnement
wecutil gs "Subscription Name"  # Afficher l'état
\end{lstlisting}

% ==============================================================================
\section{Acquisition forensique}
\label{ann:tools-acq}
% ==============================================================================

\begin{table}[H]
\centering\small
\caption{Outils d'acquisition forensique}
\label{tab:tools-acq}
\rowcolors{2}{TableauPair}{TableauImpair}
\begin{tabular}{p{3cm}p{4.5cm}p{3cm}p{2cm}p{1.5cm}}
\toprule
\tableauHeader{Outil} & \tableauHeader{Usage} &
\tableauHeader{Plateforme} & \tableauHeader{Licence} &
\tableauHeader{Format} \\
\midrule
\textbf{dd} & Copie bit à bit universelle & Linux & OS natif & RAW \\
\textbf{dcfldd} & dd avec hachage intégré & Linux & GPL & RAW \\
\textbf{FTK Imager} & Acquisition GUI + mémoire & Windows & Gratuit & E01, RAW \\
\textbf{Guymager} & Acquisition GUI Linux & Linux & GPL & E01, AFF \\
\textbf{avml} & Dump mémoire Linux (Microsoft) & Linux & MIT & LIME \\
\textbf{LiME} & Module kernel dump RAM & Linux & GPL & LIME \\
\textbf{WinPmem} & Dump mémoire Windows & Windows & Apache 2.0 & RAW \\
\textbf{Magnet RAM Capture} & Dump mémoire Windows GUI & Windows & Gratuit & RAW \\
\bottomrule
\end{tabular}
\end{table}

\begin{lstlisting}[style=StyleTerminal,
  caption={Commandes d'acquisition: référence rapide},
  label={lst:tools-acq}]
# ─── Acquisition disque avec dcfldd ──────────────────────────────────────────
dcfldd if=/dev/sdb of=/mnt/evidence/disk.dd \
  hash=sha256 hashlog=/mnt/evidence/disk.hash \
  bs=64K conv=noerror,sync sizeprobe=if
# Vérification
sha256sum /dev/sdb > /mnt/evidence/source.sha256
diff /mnt/evidence/disk.hash /mnt/evidence/source.sha256

# ─── Acquisition mémoire avec avml ───────────────────────────────────────────
sudo ./avml /mnt/usb/ram-$(hostname)-$(date +%Y%m%d-%H%M%S).lime
sha256sum /mnt/usb/ram-*.lime

# ─── Acquisition mémoire avec WinPmem ────────────────────────────────────────
winpmem_mini_x64.exe C:\evidence\ram.raw
certutil -hashfile C:\evidence\ram.raw SHA256

# ─── Montage lecture seule (Linux) ───────────────────────────────────────────
sudo losetup -r -f disk.dd && sudo losetup -a
# Identifier offset avec : mmls disk.dd
sudo mount -o ro,loop,offset=$((SECTEUR*512)) disk.dd /mnt/forensics/

# ─── Acquisition réseau : capture PCAP ───────────────────────────────────────
tcpdump -i eth0 -w /mnt/evidence/capture.pcap \
  -C 100 -G 3600 host 10.0.1.10
\end{lstlisting}

% ==============================================================================
\section{Analyse forensique}
\label{ann:tools-analysis}
% ==============================================================================

\begin{table}[H]
\centering\small
\caption{Outils d'analyse forensique}
\label{tab:tools-analysis}
\rowcolors{2}{TableauPair}{TableauImpair}
\begin{tabular}{p{3cm}p{5cm}p{2.5cm}p{2.5cm}}
\toprule
\tableauHeader{Outil} & \tableauHeader{Usage principal} &
\tableauHeader{Plateforme} & \tableauHeader{Licence} \\
\midrule
\textbf{Autopsy} & Analyse forensique complète (GUI) & Linux, Windows & Apache 2.0 \\
\textbf{The Sleuth Kit} & Analyse CLI système de fichiers & Linux & CPL \\
\textbf{Volatility~3} & Analyse mémoire vive & Linux & Volatility Lic. \\
\textbf{Plaso / log2timeline} & Timeline multi-sources & Linux & Apache 2.0 \\
\textbf{bulk\_extractor} & Extraction patterns dans image & Linux & Public Domain \\
\textbf{Scalpel} & File carving & Linux & GPL \\
\textbf{PhotoRec} & Récupération fichiers multi-format & Linux, Windows & GPL \\
\textbf{RegRipper} & Analyse registre Windows & Linux, Windows & GPL \\
\textbf{MFTECmd} & Analyse MFT NTFS & Windows & MIT \\
\textbf{PECmd} & Analyse Prefetch & Windows & MIT \\
\textbf{LECmd} & Analyse fichiers LNK & Windows & MIT \\
\textbf{SBECmd} & Analyse Shellbags & Windows & MIT \\
\textbf{EvtxECmd} & Analyse .evtx hors-ligne & Windows & MIT \\
\textbf{SIFT Workstation} & Distribution DFIR complète & Linux (VM) & SANS \\
\bottomrule
\end{tabular}
\end{table}

\begin{lstlisting}[style=StyleTerminal,
  caption={Commandes essentielles: The Sleuth Kit},
  label={lst:tools-tsk}]
# Structure des partitions
mmls image.dd

# Liste de tous les fichiers (format bodyfile pour mactime)
fls -r -m "/" image.dd > bodyfile.txt

# Timeline sur une date
mactime -b bodyfile.txt -d 2026-04-12 -z UTC > timeline.csv

# Fichiers supprimés
fls -r -d image.dd

# Extraire un fichier par inode
icat image.dd INODE_NUMBER > fichier.out

# Hash d'un fichier dans l'image
hfind -i nsrl-md5.idx hash_de_recherche
\end{lstlisting}

\begin{lstlisting}[style=StyleTerminal,
  caption={Commandes essentielles: Volatility~3},
  label={lst:tools-vol3}]
# Alias pour raccourcir
V3="python3 /opt/volatility3/vol.py"

# Windows
$V3 -f ram.mem windows.info
$V3 -f ram.mem windows.pslist
$V3 -f ram.mem windows.pstree
$V3 -f ram.mem windows.psscan         # Détecter DKOM
$V3 -f ram.mem windows.netstat
$V3 -f ram.mem windows.netscan
$V3 -f ram.mem windows.dlllist --pid PID
$V3 -f ram.mem windows.malfind        # Injection mémoire
$V3 -f ram.mem windows.hashdump       # NTLM hashes
$V3 -f ram.mem windows.cmdline        # Arguments des processus

# Linux
$V3 -f ram.lime linux.info
$V3 -f ram.lime linux.pslist
$V3 -f ram.lime linux.pstree
$V3 -f ram.lime linux.netstat
$V3 -f ram.lime linux.bash            # Historique bash en RAM
$V3 -f ram.lime linux.lsmod           # Modules kernel (rootkits)

# Comparer pslist vs psscan (DKOM)
$V3 -f ram.mem windows.pslist | awk '{print $2}' | sort > pids_list.txt
$V3 -f ram.mem windows.psscan | awk '{print $2}' | sort > pids_scan.txt
comm -13 pids_list.txt pids_scan.txt  # PIDs dans scan mais pas list
\end{lstlisting}

% ==============================================================================
\section{Analyse réseau}
\label{ann:tools-network}
% ==============================================================================

\begin{lstlisting}[style=StyleTerminal,
  caption={Commandes essentielles: Wireshark / tshark},
  label={lst:tools-wireshark}]
# ─── tshark : analyse en ligne de commande ────────────────────────────────────
# Statistiques des conversations
tshark -r capture.pcap -q -z conv,tcp | head -30

# Afficher les sessions HTTP
tshark -r capture.pcap -Y "http" \
  -T fields -e ip.src -e ip.dst \
  -e http.request.method -e http.request.full_uri

# Extraire les fichiers HTTP
tshark -r capture.pcap --export-objects http,/tmp/http_objects/

# Analyser les échanges DNS
tshark -r capture.pcap -Y "dns" \
  -T fields -e frame.time -e ip.src -e dns.qry.name -e dns.resp.addr

# Détecter les connexions longues (beacon C2 probable)
tshark -r capture.pcap -q -z conv,tcp | \
  awk '$NF > 3600'    # Durée > 1h

# ─── Zeek sur PCAP ────────────────────────────────────────────────────────────
cd /tmp/zeek_analysis/
zeek -r /evidence/capture.pcap local
ls -la   # conn.log, dns.log, http.log, ssl.log, files.log

# IPs destinations les plus fréquentes (beacon)
awk '{print $5}' conn.log | sort | uniq -c | sort -rn | head -20

# Domaines DNS avec forte entropie (DGA)
awk 'NR>1{print $10}' dns.log | \
  python3 -c "
import sys, math, collections
for d in sys.stdin:
    d=d.strip()
    if len(d)<8: continue
    c=collections.Counter(d)
    e=-sum(v/len(d)*math.log2(v/len(d)) for v in c.values())
    if e>3.5: print(f'{e:.2f}  {d}')
" | sort -rn | head -20
\end{lstlisting}

% ==============================================================================
\section{Threat Intelligence}
\label{ann:tools-cti}
% ==============================================================================

\begin{lstlisting}[style=StyleTerminal,
  caption={Commandes essentielles: MISP API et CTI},
  label={lst:tools-cti}]
# ─── Lookup d'une IP dans MISP (CLI) ─────────────────────────────────────────
curl -s -H "Authorization: $MISP_KEY" \
  -H "Accept: application/json" \
  "https://misp.lizbiz.local/attributes/restSearch/value:203.0.113.42" | \
  python3 -m json.tool

# ─── Export IoC en CSV (dernières 24h) ────────────────────────────────────────
curl -s -H "Authorization: $MISP_KEY" \
  "https://misp.lizbiz.local/attributes/restSearch/\
returnFormat:csv/type:ip-dst/last:24h/to_ids:1" -o /tmp/ioc_ips.csv

# ─── Vérifier un hash sur VirusTotal (API gratuite = 4 req/min) ──────────────
VT_KEY="votre_cle_api"
HASH="sha256_du_fichier"
curl -s -H "x-apikey: $VT_KEY" \
  "https://www.virustotal.com/api/v3/files/$HASH" | \
  python3 -c "
import sys,json
d=json.load(sys.stdin)
stats=d['data']['attributes']['last_analysis_stats']
print(f'Malicious: {stats[\"malicious\"]} / {sum(stats.values())}')
"

# ─── Abuse.ch : lookup C2 (Feodo Tracker) ────────────────────────────────────
curl -s "https://feodotracker.abuse.ch/downloads/ipblocklist.csv" | \
  grep "185.220.101.5"
\end{lstlisting}

% ==============================================================================
\section{Durcissement rapide}
\label{ann:tools-hardening}
% ==============================================================================

\begin{lstlisting}[style=StyleTerminal,
  caption={Commandes de durcissement post-incident},
  label={lst:tools-hardening}]
# ─── Réinitialiser les mots de passe (Linux) ─────────────────────────────────
# Forcer le changement au prochain login
chage -d 0 username
# Verrouiller un compte
usermod -L username
# Lister les comptes sans mot de passe (vulnérables)
awk -F: '($2 == "" || $2 == "!") {print $1}' /etc/shadow

# ─── Désactiver l'authentification par mot de passe SSH ──────────────────────
sed -i 's/^#*PasswordAuthentication.*/PasswordAuthentication no/' \
  /etc/ssh/sshd_config
sed -i 's/^#*PermitRootLogin.*/PermitRootLogin no/' \
  /etc/ssh/sshd_config
systemctl restart sshd

# ─── Désactiver les comptes de service inutilisés (Windows) ──────────────────
# Lister tous les comptes actifs
Get-ADUser -Filter {Enabled -eq $true} -Properties LastLogonDate | \
  Where {$_.LastLogonDate -lt (Get-Date).AddDays(-90)} | \
  Select Name, LastLogonDate | Sort LastLogonDate

# Désactiver en masse
Get-ADUser -Filter {Enabled -eq $true} -Properties LastLogonDate | \
  Where {$_.LastLogonDate -lt (Get-Date).AddDays(-90)} | \
  Disable-ADAccount

# ─── Activer MFA SSH (Google Authenticator) ──────────────────────────────────
apt install libpam-google-authenticator
# Ajouter dans /etc/pam.d/sshd :
# auth required pam_google_authenticator.so
# Dans /etc/ssh/sshd_config :
# ChallengeResponseAuthentication yes
# AuthenticationMethods publickey,keyboard-interactive
\end{lstlisting}
